\subsection{EXTENSION OF SCALARS FOR A VECTOR SPACE}

Recall (I, § 9, no. 1, Theorem 2) that a homomorphism of a field K into a non-zero ring A is necessarily injective.

PROPOSITION 18. Let \(\rho\) be a homomorphism of a field \(K\) into a ring \(A\) . For every exact sequence of vector \(K\) -spaces and \(K\) -linear mappings

\[
\mathrm{E} ^ {\prime} \xrightarrow {u} \mathrm{E} \xrightarrow {v} \mathrm{E} ^ {\prime \prime}
\]

the sequence

\[
\mathrm{E} _ {(\mathrm{A})} ^ {\prime} \xrightarrow {u _ {(\mathrm{A})}} \mathrm{E} _ {(\mathrm{A})} \xrightarrow {v _ {(\mathrm{A})}} \mathrm{E} _ {(\mathrm{A})} ^ {\prime \prime}
\]

is exact.

This is a special case of no. 7, Proposition 14, taking account of § 1, no. 4, Remark 4.

COROLLARY. For every K-linear mapping \(f: \mathrm{E}' \to \mathrm{E}\) , \(\operatorname{Im}(f_{(\mathbf{A})}) = (\operatorname{Im}(f))_{(\mathbf{A})}\) , \(\operatorname{Ker}(f_{(\mathbf{A})}) = (\operatorname{Ker}(f))_{(\mathbf{A})}\) , \(\operatorname{Coker}(f_{(\mathbf{A})}) = (\operatorname{Coker}(f))_{(\mathbf{A})}\) , to within canonical isomorphisms.

PROPOSITION 19. Let \(\rho\) be an injective homomorphism of a field K into a ring A. For every left vector space E over K, the canonical mapping \(\phi: \mathbf{E} \to \rho^{*}(\mathbf{E}) = \mathbf{A} \otimes_{\mathbb{K}} \mathbf{E}\) is injective. Moreover, for every vector subspace \(\mathbf{E}'\) of E, \(\rho^{*}(\mathbf{E}') = \mathbf{A} \otimes_{\mathbb{K}} \mathbf{E}'\) is canonically identified with a direct factor sub-A-module of \(\mathbf{A} \otimes_{\mathbb{K}} \mathbf{E}\) and, with this identification,

\[
(\mathbf {A} \otimes_ {\mathbf {K}} \mathbf {E} ^ {\prime}) \cap \phi (\mathbf {E}) = \phi (\mathbf {E} ^ {\prime}).\tag{35}
\]

The first assertion is a special case of § 5, no. 1, Proposition 4; the second is a special case of no. 7, Proposition 14; finally, to show (35), it suffices to take in A (considered as a right K-module) a basis \((a_{\lambda})_{\lambda \in L}\) such that \(a_{\lambda_0} = 1\) for some index \(\lambda_0\) (no. 1, Theorem 2); the elements of A \(\otimes_{\mathbb{K}}\) E can be written uniquely as \(\sum_{\lambda} a_{\lambda} \otimes x_{\lambda}\) with \(x_{\lambda} \in \mathbf{E}\) and, for such an element to belong to A \(\otimes_{\mathbb{K}}\) E', it is necessary and sufficient that \(x_{\lambda} \in \mathbf{E}'\) for all \(\lambda\) . On the other hand, the elements of \(\phi(\mathbf{E})\) are those for which \(x_{\lambda} = 0\) for \(\lambda \neq \lambda_0\) ; for an element \(\sum_{\lambda} a_{\lambda} \otimes x_{\lambda}\) to belong to (A \(\otimes_{\mathbb{K}}\) E') \(\cap\) \(\phi(\mathbf{E})\) , it is necessary and sufficient that \(x_{\lambda} = 0\) for \(\lambda \neq \lambda_0\) and \(x_{\lambda_0} \in \mathbf{E}'\) , whence the conclusion.

COROLLARY. Let \(\rho\) be an injective homomorphism of a field K into a ring A. For a K-linear mapping \(f: \mathbf{E} \to \mathbf{F}\) (where E and F are two vector spaces over K) to be injective (resp. surjective, zero), it is necessary and sufficient that \(f_{(\mathbf{A})}: \mathbf{E}_{(\mathbf{A})} \to \mathbf{F}_{(\mathbf{A})}\) be injective (resp. surjective, zero).

This follows immediately from Proposition 19 and the Corollary to Proposition 18.

PROPOSITION 20. Let \(\rho\) be a homomorphism of a field \(\mathbf{K}\) into a ring \(\mathbf{A}\) . For every left vector space \(\mathbf{E}\) over \(\mathbf{K}\) , the canonical right A-module homomorphism

\[
\upsilon : \left(\mathrm{E} ^ {*}\right) _ {(\mathrm{A})} \rightarrow \left(\mathrm{E} _ {(\mathrm{A})}\right) ^ {*}
\]

(§ 5, no. 4) is injective; it is bijective when E is finite-dimensional.

The second assertion follows from § 5, no. 4, Proposition 8. To prove the first, we note that every element of \((\mathbf{E}^{*})_{(\mathbf{A})}\) may be written as \(\sum_{i} x_{i}^{*} \otimes \alpha_{i}\) , where \(\alpha_{i} \in \mathbf{A}\) and \((x_{i}^{*})_{1 \leqslant i \leqslant n}\) is a free family in \(\mathbf{E}^{*}\) ; there corresponds to it in \((\mathbf{E}_{(\mathbf{A})})^{*}\) the linear form \(y^{*}\) such that \(y^{*}(1 \otimes x) = \sum_{i} \rho(\langle x, x_{i}^{*} \rangle)\alpha_{i}\) for all \(x \in \mathbf{E}\) . But, there exists in E a family \((x_{i})_{1 \leqslant i \leqslant n}\) such that \(\langle x_{i}, x_{j}^{*} \rangle = \delta_{ij}\) (no. 5, Corollary 2 to Theorem 7), whence \(y^{*}(1 \otimes x_{i}) = \alpha_{i}\) ; the relation \(y^{*} = 0\) therefore implies \(\alpha_{i} = 0\) for all \(i\) , which proves our assertion.

PROPOSITION 21. Let K be a field and L an extension field of K.

\begin{enumerate}
\def\labelenumi{(\roman{enumi})}
\item
  For every vector space \(\mathbf{E}\) over \(\mathbf{K}\) , \(\dim_{\mathbf{L}}(\mathrm{E}_{(\mathbf{L})}) = \dim_{\mathbf{K}}\mathbf{E}\) .
\item
  For every K-linear mapping \(u: \mathbf{E} \to \mathbf{F}\) , where \(\mathbf{E}\) and \(\mathbf{F}\) are vector spaces over \(\mathbf{K}\) , \(\operatorname{rg}(u_{(L)}) = \operatorname{rg}(u)\) .
\end{enumerate}

If \((e_{\iota})_{\iota\in I}\) is a basis of \(E\) over \(K\) , \((1\otimes e_{\iota})_{\iota\in I}\) is a basis of \(\mathbf{E}_{(L)}\) over \(L\) (§ 5, no. 1, Proposition 4), whence the first assertion; the second follows from the first and the fact that \(u_{(L)}(\mathbf{E}_{(L)})\) is canonically identified with \((u(\mathbf{E}))_{(L)}\) by the Corollary to Proposition 18.

PROPOSITION 22. Let K be a commutative field, \(\rho: \mathrm{K} \to \mathrm{A}\) an injective central homomorphism and E, F two vector spaces over K. Then the canonical homomorphism

\[
\omega : \mathrm{A} \otimes_ {\mathrm{K}} \operatorname{Hom} (\mathrm{E}, \mathrm{F}) \rightarrow \operatorname{Hom} _ {\mathrm{A}} \left(\mathrm{E} _ {(\mathrm{A})}, \mathrm{F} _ {(\mathrm{A})}\right)\tag{36}
\]

(§ 5, no. 3, formula (17)) is injective; it is bijective if A or E is a finite-dimensional vector space over K.

This is a particular case of § 5, no. 3, Proposition 7.

